\documentclass[aspectratio=169,10pt]{beamer}
\usetheme{Madrid}
\usecolortheme{seahorse}
\usefonttheme{professionalfonts}

\usepackage[utf8]{inputenc}
\usepackage[T1]{fontenc}
\usepackage{amsmath,amssymb,amsthm}
\usepackage{graphicx}
\usepackage{xcolor}
\usepackage{listings}
\usepackage{booktabs}
\usepackage{array}

\graphicspath{{figures/}}

% ---------------------------------------------------------------------
% listings style
% ---------------------------------------------------------------------
\lstset{
  language=Python,
  basicstyle=\ttfamily\scriptsize,
  numbers=left,
  numberstyle=\tiny\color{gray},
  frame=single,
  backgroundcolor=\color{gray!7},
  commentstyle=\color{green!50!black}\itshape,
  keywordstyle=\color{blue}\bfseries,
  stringstyle=\color{purple},
  showstringspaces=false,
  breaklines=true,
  tabsize=2
}

% ---------------------------------------------------------------------
% beamer cosmetics
% ---------------------------------------------------------------------
\setbeamertemplate{footline}[frame number]
\setbeamertemplate{navigation symbols}{}

% ---------------------------------------------------------------------
% title data
% ---------------------------------------------------------------------
\title[E-Values -- Ch.4]{Hypothesis Testing with E-Values}
\subtitle{Chapter 4: Post-hoc Testing and Decision Making}
\author{Based on Ramdas \& Wang}
\date{}

\begin{document}

% =======================================================================
\begin{frame}
  \titlepage
\end{frame}

\begin{frame}{Outline}
  \tableofcontents
\end{frame}

% =======================================================================
\section*{Notation and Motivation}

\begin{frame}[allowframebreaks]{Notation you'll need}
  \textbf{This chapter is about what you are allowed to do \emph{after} you have
  already looked at your data.} Four ideas recur throughout; here they are in
  plain English before we see a single formula.

  \begin{block}{Data-dependent error level $\hat\alpha$}
    Normally you fix a significance level $\alpha$ (e.g. $0.05$) \emph{before}
    collecting data. A \emph{data-dependent} level $\hat\alpha$ is a random
    variable (taking values in $(0,1)$) that is allowed to depend on the data
    in an arbitrary, even adversarial, way --- for instance, you could look at
    how convincing your result is and only then decide how strict a bar to
    hold yourself to.
  \end{block}

  \begin{block}{Post-hoc validity}
    A testing procedure is \emph{post-hoc valid} if its error-control guarantee
    survives no matter how $\hat\alpha$ was chosen, even if it was chosen
    by a mischievous statistician trying to make the result look as
    significant as possible. Analogy: a post-hoc valid guarantee is like a
    receipt that is still honoured by the shop no matter which register
    you take it to, versus an ordinary p-value which is only honoured at the
    one specific register (the one specific $\alpha$) you announced before
    walking in.
  \end{block}

  \framebreak

  \begin{block}{Optional continuation}
    You ran an experiment; the evidence was suggestive but not conclusive.
    \emph{Optional continuation} is the (data-dependent!) decision to collect
    \emph{more} data and combine it with what you already have. The key
    question: can you still control your type-I error if the decision to
    keep going --- and how much more to collect --- was itself based on
    peeking at the first batch of data? Analogy: a detective who, after an
    inconclusive round of evidence-gathering, decides (based on what they've
    found so far) whether to keep investigating, and how.
  \end{block}

  \begin{block}{Principal-agent setup}
    Two parties with different information and different incentives:
    a \emph{principal} (e.g.\ a regulator granting licenses) who sets the
    rules, and an \emph{agent} (e.g.\ a pharmaceutical company) who reacts to
    them, possibly by lying about what they believe or by running trials they
    know are futile. \emph{Incentive compatibility} means the rules are
    designed so that agents who know their claim is false have no reason to
    even try.
  \end{block}

  \textbf{Recall from earlier chapters:} an \emph{e-variable} $E$ for a null
  class $\mathcal P$ satisfies $\mathbb{E}^{\mathbb P}[E]\le 1$ for every
  $\mathbb P\in\mathcal P$; a \emph{p-variable} $P$ satisfies
  $\mathbb{P}(P\le \alpha)\le\alpha$ for every $\alpha\in(0,1)$ and every
  $\mathbb P\in\mathcal P$ (equivalently $\mathbb E^{\mathbb P}[\mathbb 1\{P\le\alpha\}]\le\alpha$).
\end{frame}

\begin{frame}{Why does this chapter matter?}
  \begin{itemize}
    \item Standard hypothesis testing textbook theory assumes: fix $\alpha$,
    fix your test, fix your sample size, collect data \emph{once}, decide.
    \item Real science never works this way: people peek at results,
    change thresholds, run follow-up experiments only when the first one
    looks promising, and make more than one decision from a single dataset.
    \item \textbf{Main message of Chapter 4:} e-values are exactly the
    right currency for all of this --- they remain valid under
    data-dependent significance levels, non-binary decisions, data-dependent
    decision problems, optional continuation, and even when the incentives
    of the analyst are adversarial. P-values, in general, are \emph{not}.
    \item We will make every one of these claims precise, and back each one
    with a worked numerical example on a single running scenario: testing
    whether a coin is fair, $H_0:\theta=0.5$.
  \end{itemize}
\end{frame}

% =======================================================================
\section{Testing at Data-Dependent Error Levels}

\begin{frame}{Section 4.1}
  \tableofcontents[currentsection]
\end{frame}

\begin{frame}{Intuition: what if $\alpha$ is chosen after seeing the data?}
  \begin{itemize}
    \item Recall: a level-$\alpha$ test for $\mathcal P$ is a function
    $\phi$ (taking values in $\{0,1\}$, or more generally $[0,1]$, the
    probability of rejecting) such that $\mathbb E^{\mathbb P}[\phi]\le\alpha$
    for every $\mathbb P\in\mathcal P$.
    \item As presented in earlier chapters, $\alpha$ was always a
    \emph{predefined constant}, agreed upon before the experiment.
    \item \textbf{Question:} can we build a family of tests
    $(\phi_\alpha)_{\alpha\in(0,1)}$, one for every possible level, such that
    we're allowed to pick \emph{which} $\alpha$ to report --- using a rule
    $\hat\alpha$ that can depend on the data in \emph{any} way --- and still
    have $\mathbb E^{\mathbb P}[\phi_{\hat\alpha}]\le\hat\alpha$ hold?
    \item This is a genuinely stronger requirement than ordinary level-$\alpha$
    validity, because $\hat\alpha$ is allowed to depend on $\phi_\alpha$ itself
    for every $\alpha$ simultaneously (i.e.\ on the whole family, not just one
    member).
  \end{itemize}
\end{frame}

\begin{frame}[allowframebreaks]{Definition 4.1: post-hoc valid tests and p-variables}
  \begin{block}{Post-hoc valid family of tests}
    A collection $(\phi_\alpha)_{\alpha\in(0,1)}$ is a \emph{post-hoc valid
    family of tests} for $\mathcal P$ if for every $\mathbb P\in\mathcal P$
    and every random variable $\hat\alpha$ taking values in $(0,1)$,
    \[
      \mathbb E^{\mathbb P}\!\left[\frac{\phi_{\hat\alpha}}{\hat\alpha}\right]\le 1.
    \]
    Equivalently,
    \[
      \mathbb E^{\mathbb P}\!\left[\sup_{\alpha\in(0,1)} \frac{\phi_\alpha}{\alpha}\right]\le 1
      \qquad\text{for all }\mathbb P\in\mathcal P.
    \]
    \emph{Symbols:} $\phi_\alpha$ is the test at nominal level $\alpha$
    (rejection probability); $\hat\alpha$ is any --- possibly adversarial
    --- data-dependent level; the sup is over all possible levels at once.
  \end{block}

  \begin{block}{Post-hoc p-variable}
    A nonnegative random variable $P$ is a \emph{post-hoc p-variable} for
    $\mathcal P$ if for every $\mathbb P\in\mathcal P$ and every random
    variable $\hat\alpha\in(0,1)$,
    \[
      \mathbb E^{\mathbb P}\!\left[\frac{\mathbb 1\{P\le\hat\alpha\}}{\hat\alpha}\right]\le 1,
      \qquad\text{equivalently}\qquad
      \mathbb E^{\mathbb P}\!\left[\sup_{\alpha\in(0,1)}\frac{\mathbb 1\{P\le\alpha\}}{\alpha}\right]\le 1.
      \tag{4.1}
    \]
  \end{block}

  \framebreak

  \textbf{Contrast with an ordinary p-variable.} For a p-variable $P$, only a
  \emph{weaker} condition holds:
  \[
    \sup_{\alpha\in(0,1)} \mathbb E^{\mathbb P}\!\left[\frac{\mathbb 1\{P\le\alpha\}}{\alpha}\right]\le 1
    \qquad\text{for all }\mathbb P\in\mathcal P. \tag{4.2}
  \]
  This is precisely the usual definition of a p-variable. The crucial
  difference between (4.1) and (4.2) is whether the $\sup_\alpha$ sits
  \emph{inside} or \emph{outside} the expectation:
  \begin{itemize}
    \item (4.2): for \emph{each fixed} $\alpha$ chosen in advance, the average
    rejection rate is $\le\alpha$. Fine if you commit to $\alpha$ up front.
    \item (4.1): even the \emph{worst-case, data-dependent} choice of
    $\alpha$ cannot inflate the average past $1$. This is what post-hoc
    validity requires, and it is a fundamentally stronger property.
  \end{itemize}

  \begin{block}{Remark 4.2 (a useful generalization)}
    For a p-variable $P$, the condition $\mathbb E^{\mathbb P}[\mathbb 1\{P\le T\}/T]\le 1$
    (with $0/0:=0$) also holds for some nonnegative random variables $T$ (not
    necessarily in $(0,1)$), provided
    \[
      p\mapsto \mathbb P(T\le t\mid P\le p)\ \text{is increasing for all } t\in\mathbb R. \tag{4.4}
    \]
    This is a form of \emph{negative dependence} between $P$ and $T$.
    Examples satisfying (4.4): (a) $T$ is a constant; (b) $T$ is independent
    of $P$; (c) $T$ is a decreasing function of $P$.
  \end{block}
\end{frame}

\begin{frame}[allowframebreaks]{Proposition 4.3: e-variables generate post-hoc valid families}
  \begin{block}{Proposition 4.3}
    Let $E$ be an e-variable for $\mathcal P$. Then $((\alpha E)\wedge 1)_{\alpha\in(0,1)}$
    is a post-hoc valid family of tests, and $(\mathbb 1\{E\ge 1/\alpha\})_{\alpha\in(0,1)}$
    is a post-hoc valid family of binary tests.

    \smallskip
    Conversely, every post-hoc valid family of tests is dominated by
    $((\alpha E)\wedge 1)_{\alpha\in(0,1)}$ for some e-variable $E$, and every
    post-hoc valid family of binary tests is dominated by
    $(\mathbb 1\{E\ge 1/\alpha\})_{\alpha\in(0,1)}$ for some e-variable $E$.
  \end{block}
  \emph{Symbols:} $x\wedge 1 = \min(x,1)$; ``dominated'' means member-wise
  $\le$ (a dominated family rejects no more often, hence is weakly less
  useful). \textbf{Punchline:} e-variables are (essentially) the
  \emph{only} way to build post-hoc valid families.

  \framebreak

  \textbf{Why is this true? (proof idea)} The key elementary fact used
  throughout this chapter is
  \[
    \mathbb 1\{x\ge 1\} \le x\wedge 1 \le x \qquad \text{for } x\ge 0, \tag{4.5}
  \]
  which is exactly the inequality underlying Markov's inequality. Applying
  (4.5) with $x = E/\hat\alpha$, for any random $\hat\alpha\in(0,1)$:
  \[
    \mathbb E^{\mathbb P}\!\left[\frac{\mathbb 1\{E\ge 1/\hat\alpha\}}{\hat\alpha}\right]
    \le \mathbb E^{\mathbb P}\!\left[\frac{(\hat\alpha E)\wedge 1}{\hat\alpha}\right]
    \le \mathbb E^{\mathbb P}\!\left[\frac{\hat\alpha E}{\hat\alpha}\right]
    = \mathbb E^{\mathbb P}[E] \le 1.
  \]
  This one chain of inequalities is the engine behind almost every result in
  this section: it shows both families are post-hoc valid \emph{for any}
  $\hat\alpha$ whatsoever, adversarial or not, because $\hat\alpha$ simply
  cancels out.

  \textit{Converse (sketch):} given a post-hoc valid family
  $(\phi_\alpha)$, set $E=\sup_\alpha \phi_\alpha/\alpha$; post-hoc validity
  says exactly that $E$ is an e-variable, and by construction
  $\phi_\alpha\le(\alpha E)\wedge 1$.
\end{frame}

\begin{frame}{Proposition 4.4: post-hoc p-variables are (essentially) $1/E$}
  \begin{block}{Proposition 4.4}
    For a random variable $P$, the following are equivalent:
    \begin{enumerate}
      \item $P$ is a post-hoc p-variable for $\mathcal P$;
      \item $P = 1/E$ on the event $\{P<1\}$, for some e-variable $E$ for $\mathcal P$;
      \item $P \ge (1/E)\wedge 1$ for some e-variable $E$ for $\mathcal P$.
    \end{enumerate}
  \end{block}
  If instead we strengthen (4.1) by letting $\alpha$ range over $(0,\infty)$,
  \[
    \sup_{\alpha\in(0,\infty)} \mathbb E^{\mathbb P}\!\left[\frac{\mathbb 1\{P\le\alpha\}}{\alpha}\right]\le 1
    \quad\text{for all }\mathbb P\in\mathcal P, \tag{4.6}
  \]
  then a post-hoc p-variable is \emph{precisely} $P=1/E$ for some e-variable
  $E$ (which may exceed $1$, since we allow p-variables larger than $1$).
  \textbf{Takeaway:} post-hoc p-variables and e-variables are really the same
  object, viewed from two sides of the same coin: $P = 1/E$.
\end{frame}

\begin{frame}{Interpreting the stronger guarantee: $E=41$ vs.\ $P=0.024$}
  Suppose on the same hypothesis we observe an e-value $E\approx 41.05$
  \emph{or} a p-value $P\approx 0.0244$ on the same data (as we compute
  concretely below). These look ``comparable'' ($1/E\approx P$), but they are
  \emph{not} interchangeable:
  \begin{itemize}
    \item Because of Proposition 4.3--4.4, we can \emph{safely} say: ``we
    reject $H_0$ at level $0.0244$ (chosen after seeing $E$!), and even at
    level $0.05$.'' This claim is protected by post-hoc validity.
    \item We \emph{cannot} safely say: ``we reject $H_0$ at level $0.05$''
    based on $P\approx 0.0244$, \emph{unless the level $0.05$ was fixed
    before collecting and inspecting the data.}
    \item If the level was chosen only \emph{after} peeking (e.g.\ ``the
    p-value looked good, so let's call $0.05$ our threshold''), the ordinary
    p-value's guarantee (4.2) says nothing, because (4.2) never promised
    anything about a $\hat\alpha$ that depends on the data.
  \end{itemize}
  \textbf{Rule of thumb:} an e-value of $1/a$ carries \emph{strictly more}
  information than a p-value of $a$, precisely because of this extra
  post-hoc guarantee.
\end{frame}

\begin{frame}{Running example, Part 1: peeking is safe with e-values}
  \textbf{Setup.} $H_0:\theta=0.5$ (fair coin). A scientist tosses the coin
  $n=20$ times and forms the likelihood-ratio e-value against the simple
  alternative $\theta=0.75$:
  \[
    E(x) = \left(\frac{0.75}{0.5}\right)^{H}\left(\frac{0.25}{0.5}\right)^{T} = 1.5^{H}\cdot 0.5^{T},
  \]
  where $H$ = number of heads, $T=20-H$ = number of tails. This $E$ is an
  e-variable for $H_0$: $\mathbb E^{\theta=0.5}[E] = 1$ exactly (it's a
  likelihood ratio).

  \medskip
  \textbf{Data.} Suppose she observes $H=16$ heads, $T=4$ tails. Then
  \[
    E = 1.5^{16}\times 0.5^{4} \approx 41.05.
  \]
  \textbf{Post-hoc decision.} \emph{After} seeing $E\approx 41.05$, she
  decides --- on the spot, informed by how good the number looks --- to
  report significance at the data-dependent level
  \[
    \hat\alpha = 1/E \approx 0.0244
  \]
  (the smallest level her evidence supports). By Proposition 4.3--4.4 this is
  completely valid: the family $(\mathbb 1\{E\ge1/\alpha\})_\alpha$ is
  post-hoc valid \emph{for any} rule choosing $\hat\alpha$, including this
  self-serving one. She may honestly report ``significant at level $0.0244$,
  chosen after seeing the data.''
\end{frame}

\begin{frame}{Running example, Part 2: peeking breaks p-value validity}
  \textbf{Same idea, now with an ordinary p-value.} Under $H_0$, suppose $P$
  is (as usual) continuous with $P\sim\text{Uniform}(0,1)$.

  \medskip
  \textbf{The adversarial choice.} Consider a scientist who, having computed
  $P$, sets her ``after-the-fact'' threshold to be exactly her own p-value:
  $\hat\alpha := P$. Then trivially
  \[
    \mathbb P\big(P \le \hat\alpha\big) = \mathbb P(P\le P) = 1,
  \]
  i.e.\ she can \emph{always} claim ``significant at level $\hat\alpha$'' ---
  a $100\%$ false-rejection rate under the null, however small $\alpha$
  ``looks''! Compare to a pre-registered $\alpha=0.05$: type-I error is
  exactly $0.05$, as it should be.

  \medskip
  \textbf{Same fact, via Definition 4.1.} The quantity that post-hoc validity
  requires to be $\le 1$ blows up for p-values:
  \[
    \mathbb E^{\mathbb P}\!\left[\sup_{\alpha\in(0,1)}\frac{\mathbb 1\{P\le\alpha\}}{\alpha}\right]
     = \mathbb E^{\mathbb P}\!\left[\frac1P\right]
     = \int_0^1 \frac1p\,dp = \infty.
  \]
  \textbf{Contrast:} for the e-value family, the analogous quantity is
  $\mathbb E^{\mathbb P}[E\cdot\mathbb 1\{E\ge1\}]\le\mathbb E^{\mathbb P}[E]\le 1$ ---
  finite, by Markov's inequality, no matter what. This is the concrete,
  numerical face of ``post-hoc validity for e-values, but not for p-values.''
\end{frame}

\begin{frame}[fragile]{Python: post-hoc risk -- e-values vs.\ p-values (\S4.1)}
\begin{lstlisting}
import numpy as np
from scipy.special import erfc

rng = np.random.default_rng(0)
theta0, theta1, n, M = 0.5, 0.6, 30, 200_000

# simulate under H0: theta = 0.5
heads = rng.binomial(n, theta0, size=M)
tails = n - heads

# likelihood-ratio e-value for theta1 vs theta0
E = (theta1/theta0)**heads * ((1-theta1)/(1-theta0))**tails

# continuous one-sided p-value (normal approximation)
z = (heads - n*theta0) / np.sqrt(n*theta0*(1-theta0))
P = np.clip(0.5*erfc(z/np.sqrt(2)), 1e-300, 1.0)

R_P = 1.0 / P                 # sup_alpha 1{P<=alpha}/alpha
R_E = E * (E >= 1.0)          # sup_alpha 1{E>=1/alpha}/alpha

print("running avg R_P (p-values):", np.mean(R_P))   # unstable, >> 1
print("running avg R_E (e-values):", np.mean(R_E))   # stays <= 1
\end{lstlisting}
  \textbf{Output (M=200{,}000 replicates):} running average of $R_P\approx 15$
  (unstable, keeps growing with $M$ since $\mathbb E[1/P]=\infty$), while
  running average of $R_E\approx 0.71 \le 1$, exactly as guaranteed by
  Markov's inequality applied to the e-variable $E$. See
  \texttt{figures/fig\_posthoc\_validity.pdf} for the full trajectory.
\end{frame}

\begin{frame}{Section 4.1 take-aways}
  \begin{itemize}
    \item Ordinary level-$\alpha$ tests only control error \emph{for a fixed,
    pre-registered} $\alpha$.
    \item Post-hoc validity (Def.\ 4.1) additionally controls error for
    \emph{any}, even adversarially data-dependent, choice of level $\hat\alpha$.
    \item e-variables generate post-hoc valid families essentially uniquely
    (Prop.\ 4.3), and post-hoc p-variables are just $1/E$ (Prop.\ 4.4).
    \item Concretely: it is fine to look at your e-value and then decide how
    impressive a claim to make; doing the analogous thing with an ordinary
    p-value can inflate your true error rate all the way to $100\%$.
  \end{itemize}
\end{frame}

% =======================================================================
\section{Post-hoc Decisions Using P-values and E-values}

\begin{frame}{Section 4.2}
  \tableofcontents[currentsection]
\end{frame}

\begin{frame}{Intuition: beyond accept/reject}
  \begin{itemize}
    \item So far, ``decisions'' meant a binary choice: reject $H_0$ or not.
    \item Real decision-making is richer: a doctor may choose among
    \emph{several} treatment intensities; a regulator may choose among
    several licensing tiers; and \emph{which} decisions are even on the
    table may itself depend on what the data show (a data-dependent
    ``scenario'').
    \item Section 4.2 asks: if we build our decision rule from a p-value vs.\
    from an e-value, which one keeps its risk-control promise once decisions
    become non-binary, or once the decision problem itself depends on data?
    \item \textbf{Spoiler:} for the simplest case (binary decisions, fixed
    scenario), both work. As soon as either non-binary decisions or
    data-dependent scenarios enter, only the e-value-based rule survives.
  \end{itemize}
\end{frame}

\begin{frame}[allowframebreaks]{Setup: loss functions and the risk-constrained problem}
  Consider loss functions $L_b:\{0,1\}\times D_b\to\mathbb R$ indexed by a
  \emph{scenario} $b\in\mathcal B\subseteq\mathbb R$, where $D_b\subseteq\mathbb N_0$
  is the set of possible decisions in scenario $b$.
  \begin{itemize}
    \item $L_b(0,d)$: loss to the decision maker from decision $d$
    \emph{when $H_0$ is true};
    \item $L_b(1,d)$: loss from decision $d$ \emph{when the alternative is true};
    \item WLOG, $d\mapsto L_b(0,d)$ is strictly increasing and
    $d\mapsto L_b(1,d)$ is strictly decreasing --- larger $d$ = a ``more
    aggressive'' decision: costlier if $H_0$ true, more beneficial if the
    alternative is true.
  \end{itemize}
  Data $X\in\mathcal Z$; decision rule $\delta(X)\in D_b$.

  \begin{block}{Risk-constrained decision problem (4.7)}
    \[
      \text{minimize}_{\delta:\mathcal Z\to D_b}\ \ \mathbb E^{\mathbb Q}[L_b(1,\delta(X))]
      \quad\text{s.t.}\quad
      \sup_{\mathbb P\in\mathcal P}\mathbb E^{\mathbb P}[L_b(0,\delta(X))]\le\gamma,
    \]
  \end{block}
  where $\gamma>0$ is a prespecified \emph{risk budget}, and the constraint
  is the \emph{risk constraint}. (We fix a simple alternative
  $\mathcal Q=\{\mathbb Q\}$ for simplicity; the composite case is analogous.)
\end{frame}

\begin{frame}{Example 4.5: ternary significance decisions}
  A very common real practice: label a result ``non-significant (NS)'',
  ``significant (S)'', or ``very significant (VS)'' after seeing the data.
  \[
    D_b=\{0,1,2\},\qquad 0=\text{NS},\ 1=\text{S},\ 2=\text{VS}.
  \]
  If $H_0$ true, decision $0$ (NS) has least loss; if the alternative is
  true, decision $2$ (VS) has least loss. A common choice:
  \[
    \gamma=1,\qquad L_b(0,0)=0,\qquad L_b(0,1)=\frac1{\alpha_1},\qquad L_b(0,2)=\frac1{\alpha_2},
  \]
  for some $0<\alpha_2<\alpha_1<1$ (e.g.\ $\alpha_1=0.05$, $\alpha_2=0.01$).
\end{frame}

\begin{frame}{Decision rules based on p-values and e-values}
  Suppose a p-variable $P(X)$ and an e-variable $E(X)$ for $\mathcal P$ are
  both computed from the data.
  \begin{block}{p-value decision rule (4.8)}
    \[
      \delta_P(X) = \max\{\,d\in D_b : P(X)\,L_b(0,d)\le\gamma\,\}.
    \]
  \end{block}
  \begin{block}{e-value decision rule (4.9)}
    \[
      \delta_E(X) = \max\{\,d\in D_b : L_b(0,d)\le\gamma\,E(X)\,\}.
    \]
  \end{block}
  \emph{Symbols:} both rules pick the most aggressive decision $d$ still
  compatible with the risk budget $\gamma$, using either $P(X)$ or $E(X)$ as
  the ``evidence multiplier.''
\end{frame}

\begin{frame}{Example 4.6: ternary decisions, worked out with numbers}
  Take $\alpha_1=0.05$, $\alpha_2=0.01$, $\gamma=1$ (as in Example 4.5).
  \begin{block}{p-value rule (4.8)}
    \[
      \delta_P(X) =
      \begin{cases}
        2 & \text{if } P(X)\le 0.01,\\
        1 & \text{if } 0.01 < P(X) \le 0.05,\\
        0 & \text{otherwise.}
      \end{cases}
    \]
    This is standard scientific practice (``$*$'' vs.\ ``$**$'' asterisks on
    p-values) -- \textbf{but it will turn out not to satisfy the risk
    constraint (4.7) in general (Example 4.10)!}
  \end{block}
  \begin{block}{e-value rule (4.9)}
    \[
      \delta_E(X) =
      \begin{cases}
        2 & \text{if } E(X)\ge 100,\\
        1 & \text{if } 20\le E(X) < 100,\\
        0 & \text{otherwise.}
      \end{cases}
    \]
    This rule \emph{always} satisfies the risk constraint (Theorem 4.9).
  \end{block}
\end{frame}

\begin{frame}[allowframebreaks]{Binary decisions in a fixed scenario: Proposition 4.7}
  Simplest case: $D_b=\{0,1\}$ ($0$=accept, $1$=reject), and WLOG
  $L_b(0,0)=L_b(1,1)=0$, $L_b(0,1)=L_b(1,0)=1$. Then (4.7) becomes the
  familiar
  \[
    \text{maximize}\ \ \mathbb Q(\delta(X)=1)\quad\text{s.t.}\quad
    \sup_{\mathbb P\in\mathcal P}\mathbb P(\delta(X)=1)\le\gamma, \tag{4.10}
  \]
  i.e.\ maximize power subject to type-I error control $\gamma$.
  \begin{itemize}
    \item p-value rule: $\delta_P(X)=1$ iff $P(X)\le\gamma$.
    \item e-value rule: $\delta_E(X)=1$ iff $E(X)\ge 1/\gamma$.
  \end{itemize}

  \begin{block}{Proposition 4.7}
    Given a fixed scenario $b$ with $D_b=\{0,1\}$, both $\delta_P$ and
    $\delta_E$ satisfy the risk constraint in (4.10).
  \end{block}

  \framebreak
  \textbf{Proof.} For any $\mathbb P\in\mathcal P$:
  \[
    \mathbb E^{\mathbb P}[L_b(0,\delta_P(X))] = \mathbb P(\delta_P(X)=1) = \mathbb P(P(X)\le\gamma)\le\gamma,
  \]
  using the p-variable property. And
  \[
    \mathbb E^{\mathbb P}[L_b(0,\delta_E(X))] \le \mathbb E^{\mathbb P}[\gamma E(X)] \le \gamma,
  \]
  using $L_b(0,\delta_E(X))\le\gamma E(X)$ by construction, and
  $\mathbb E^{\mathbb P}[E]\le1$.

  \medskip
  \textbf{Conclusion for this simple case:} both approaches are valid! Since
  $1/E(X)$ is itself a p-variable (Section 2.1), the e-value rule is
  typically \emph{more conservative} than the p-value rule here. The real
  difference emerges only once we leave this simplest binary/fixed-scenario
  world.
\end{frame}

\begin{frame}{Non-binary or data-dependent scenarios: Example 4.8}
  \textbf{Vaccine story.} A clinical trial's data can put the scientist in
  one of two \emph{scenarios} $B\in\{1,2\}$:
  \begin{itemize}
    \item $B=1$ (``promising''): decisions $D_1=\{0,1\}$, where
    $\delta(X,1)=0$ means \emph{don't vaccinate anyone} and $\delta(X,1)=1$
    means \emph{vaccinate the highest-risk population}.
    \item $B=2$ (``extremely effective'', unexpected): now
    $\delta(X,2)=1$ means something more aggressive: \emph{vaccinate the
    full population}.
  \end{itemize}
  \begin{block}{Toy numeric loss table}
    \begin{center}
    \begin{tabular}{lccc}
      \toprule
      Scenario $b$ & $L_b(0,0)$ & $L_b(0,1)$ & meaning of $d=1$\\
      \midrule
      $B=1$ (promising) & $0$ & $1$ & vaccinate high-risk only\\
      $B=2$ (extremely effective) & $0$ & $10$ & vaccinate everyone\\
      \bottomrule
    \end{tabular}
    \end{center}
  \end{block}
  Note $D_1$ and $D_2$ carry \emph{different meanings} for the same label
  $d=1$ -- and which scenario $B$ you're even in is itself random,
  determined by the data!
\end{frame}

\begin{frame}[allowframebreaks]{The formal data-dependent problems (4.11)-(4.12)}
  Write $B$ (a function of $X$) for the observable, data-dependent scenario;
  the decision rule now takes both $X$ and $B$: $\delta(X,B)\in D_B$.

  \begin{block}{Problem (4.11): scenario $B$ known to the decision maker}
    \[
      \text{minimize}\ \ \mathbb E^{\mathbb Q}[L_B(1,\delta(X,B))]
      \quad\text{s.t.}\quad
      \sup_{\mathbb P\in\mathcal P}\mathbb E^{\mathbb P}[L_B(0,\delta(X,B))]\le\gamma.
    \]
  \end{block}

  \begin{block}{Problem (4.12): worst-case objective/constraint over all $b\in\mathcal B$}
    \[
      \text{minimize}\ \ \mathbb E^{\mathbb Q}\!\Big[\sup_{b\in\mathcal B}L_b(1,\delta(X,b))\Big]
      \quad\text{s.t.}\quad
      \sup_{\mathbb P\in\mathcal P}\mathbb E^{\mathbb P}\!\Big[\sup_{b\in\mathcal B}L_b(0,\delta(X,b))\Big]\le\gamma.
    \]
  \end{block}
  The e-value rule extends naturally to this setting and by construction
  satisfies, for every $b$,
  \[
    L_b(0,\delta_E(X,b)) \le \gamma E(X). \tag{4.13}
  \]

  \framebreak
  \begin{block}{Theorem 4.9}
    \begin{enumerate}
      \item $\delta_E$ (for any e-variable $E(X)$) satisfies the risk
      constraints in both (4.11) and (4.12).
      \item Any $\delta^*$ satisfying the constraint in (4.11) can be written
      as $\delta_E(X,B)$ for some e-variable $E(X)$.
      \item Any $\delta^*$ satisfying the constraint in (4.12) satisfies
      $\delta^*(X,b)\le\delta_E(X,b)$ for some e-variable $E(X)$ and
      \emph{all} $b\in\mathcal B$ (i.e.\ $\delta^*$ is weakly improved by
      using e-variables).
      \item The p-value rule $\delta_P$ does \emph{not} satisfy the risk
      constraint in (4.11) or (4.12) in general, once $D_b$ has at least $3$
      elements or $\mathcal B$ has at least $2$ elements.
    \end{enumerate}
  \end{block}
  \textbf{Proof of (i):} directly from (4.13), $\mathbb E^{\mathbb P}[\sup_b L_b(0,\delta_E(X,b))]\le\mathbb E^{\mathbb P}[\gamma E(X)]\le\gamma$.
\end{frame}

\begin{frame}{Example 4.10: a concrete counterexample for $\delta_P$}
  Let $P(X)\sim\text{Uniform}[0,1]$ under $\mathbb P\in\mathcal P$.
  Take $\mathcal B=\{1,2\}$, $D_b=\{0,1\}$, $L_1(0,1)=\ell_1$,
  $L_2(0,1)=\ell_2$ with $\ell_1<\ell_2$, and let the scenario be
  $B=2$ exactly when $P(X)\le\gamma/\ell_2$ (otherwise $B=1$). A direct
  computation gives
  \[
    \mathbb E^{\mathbb P}[L_B(0,\delta_P(X))]
    = \ell_1\,\mathbb P\!\left(\tfrac{\gamma}{\ell_2}\le P(X)\le\tfrac{\gamma}{\ell_1}\right)
      + \ell_2\,\mathbb P\!\left(P(X)\le\tfrac{\gamma}{\ell_2}\right)
    = \gamma\Big(1-\tfrac{\ell_1}{\ell_2}\Big)+\gamma \;>\;\gamma. \tag{4.14}
  \]
  \textbf{Concrete numbers:} take $\gamma=0.05$, $\ell_1=1$, $\ell_2=10$.
  Then the achieved risk is
  \[
    0.05\Big(1-\tfrac1{10}\Big)+0.05 = 0.045+0.05 = 0.095,
  \]
  \emph{nearly double} the allowed budget $\gamma=0.05$! The p-value rule
  looks fine \emph{within} each fixed scenario, but silently breaks the
  overall risk guarantee once the scenario itself is chosen using the data.
  By contrast (Theorem 4.9(i)), $\delta_E$ never does this.
\end{frame}

\begin{frame}[fragile]{Python: verifying Example 4.10 by simulation (\S4.2)}
\begin{lstlisting}
import numpy as np
rng = np.random.default_rng(1)

gamma, ell1, ell2 = 0.05, 1.0, 10.0
M = 2_000_000
P = rng.uniform(0, 1, size=M)          # P-variable under H0

# scenario B depends on the data:
B = np.where(P <= gamma/ell2, 2, 1)

# delta_P(X) = 1 iff P(X) * L_B(0,1) <= gamma
L_B_of_1 = np.where(B == 2, ell2, ell1)
delta_P = (P * L_B_of_1 <= gamma).astype(int)

loss = np.where(delta_P == 1, L_B_of_1, 0.0)
print("empirical risk of delta_P:", loss.mean(), " (budget =", gamma, ")")
# --> approx 0.095, almost double the risk budget gamma = 0.05
\end{lstlisting}
  The simulated risk matches the closed-form value $0.095$ from (4.14),
  confirming that a p-value-based ternary/scenario-dependent rule can
  silently blow through its risk budget.
\end{frame}

\begin{frame}{Section 4.2 take-aways}
  \begin{itemize}
    \item With \emph{binary} decisions and a \emph{fixed} scenario, both
    p-value- and e-value-based decision rules are valid (Prop.\ 4.7);
    e-value rules tend to be more conservative here.
    \item As soon as decisions are non-binary (e.g.\ NS/S/VS) \emph{or} the
    decision problem itself is data-dependent (the scenario $B$ depends on
    $X$), only e-value-based rules keep the risk-control promise
    (Theorem 4.9); p-value-based rules can silently violate it
    (Example 4.10).
    \item Practical warning: this failure mode is common and often
    \emph{unconscious} -- e.g.\ marking p-values with significance asterisks
    after the fact is exactly this ternary-decision setting.
  \end{itemize}
\end{frame}

% =======================================================================
\section{Optional Continuation of Experiments}

\begin{frame}{Section 4.3}
  \tableofcontents[currentsection]
\end{frame}

\begin{frame}{Intuition: ``I should have collected more data''}
  \begin{itemize}
    \item A scientist runs a carefully pre-planned experiment. The result is
    close, but not quite significant. In hindsight, they realize they
    underestimated the effect size and should have collected more data.
    \item With a p-value, this realization is \emph{useless}: the type-I
    error budget for that test was already spent the moment they committed
    to that sample size and threshold. Tacking on more data and re-testing
    (even honestly) inflates the false-positive rate.
    \item \textbf{The question of Section 4.3:} is there a principled way to
    let a scientist decide --- \emph{after seeing the first batch of data}
    --- whether to run more, and if so, keep a valid guarantee against
    $H_0$? We call this \emph{optional continuation}.
    \item The answer, again, is yes with e-values.
  \end{itemize}
\end{frame}

\begin{frame}{Coin-fairness example: the Neyman--Pearson baseline}
  \textbf{Setup.} Test $H_0:\theta=0.5$ (fair coin) vs.\ an anticipated bias
  $\theta=0.6$. To guarantee both type-I and type-II error at $0.05$, the
  optimal (Neyman--Pearson) likelihood-ratio test:
  \begin{itemize}
    \item tosses the coin $n=280$ times;
    \item rejects $H_0$ iff at least $154$ heads appear.
  \end{itemize}
  Suppose, on performing those $280$ tosses, the scientist does \emph{not}
  reject. Maybe it's a type-II error; maybe they should have planned for a
  smaller bias, say $\theta=0.55$. \textbf{But it's too late}: the type-I
  error budget of this test has already been fully spent, and re-examining
  the same data with a new threshold is not valid. With p-values, such
  regrets are simply futile.
\end{frame}

\begin{frame}[allowframebreaks]{E-values to the rescue: combining fresh batches}
  Rewrite the original test as: reject when the LR-test p-value drops below
  $0.05$. Now instead summarize the very same $280$-toss experiment as an
  e-value $E_1$ from the start (the likelihood ratio is a natural choice).

  \begin{block}{Key martingale property}
    Suppose the scientist (or someone else) is permitted to toss the coin
    more times -- say $50$ more -- and computes a fresh e-value $E_2$ based
    \emph{only} on this new batch. Under the null,
    \[
      \mathbb E^{\mathbb P}[E_2\mid E_1] = 1.
    \]
    Then $E^{(2)} := E_2 E_1$ is \emph{also} a valid e-value for the same
    null (by the law of iterated expectation:
    $\mathbb E^{\mathbb P}[E^{(2)}]=\mathbb E^{\mathbb P}[\mathbb E^{\mathbb P}[E_2E_1\mid E_1]]=\mathbb E^{\mathbb P}[E_1\cdot1]=1$).
    Crucially: \emph{both} the decision to run this second experiment
    \emph{and} its design (e.g.\ how many more tosses) are allowed to depend
    on the first batch of data -- this is a post-hoc decision, and it
    remains valid!
  \end{block}

  \framebreak
  The scientist may still be unsatisfied. They can collect a third batch,
  summarize it as $E_3$, and form $E^{(3)} := E_3 E^{(2)}$, itself a valid
  e-value. \textbf{This can be repeated indefinitely}: the evidence process
  $E^{(1)}, E^{(2)}, \dots$ can be monitored continually, and $H_0$ rejected
  as soon as the current evidence exceeds $1/\alpha$ -- this is controlled
  by \emph{Ville's inequality} (Chapter 7). Alternatively, the scientist may
  stop at any data-dependent stopping time $\tau$ and report $E^{(\tau)}$,
  itself an e-value, by the \emph{optional stopping theorem} (Chapter 7).
  When each batch shrinks to a single toss, this becomes a fully sequential
  test -- Wald's SPRT (Section 7.2).
\end{frame}

\begin{frame}{Worked numbers: from ``promising'' to a valid rejection}
  \textbf{Continuing our coin example.} Suppose the scientist's original
  $n=280$-toss batch yields e-value
  \[
    E_1 = 12.
  \]
  Since we need $E\ge 1/\alpha = 20$ to reject at $\alpha=0.05$, they cannot
  reject yet -- but the post-hoc p-value $1/E_1\approx 0.083$ is
  ``promising.'' \emph{Because} it looked promising, they decide (post-hoc!)
  to run $50$ more tosses, obtaining a fresh e-value on the new batch alone:
  \[
    E_2 = 2.5.
  \]
  The combined e-value is
  \[
    E^{(2)} = E_1\times E_2 = 12\times 2.5 = 30 \ge 20,
  \]
  so \textbf{they may now validly reject $H_0$ at level $\alpha=0.05$} --
  even though the second experiment was only run \emph{because} the first
  one was almost significant, and even though this is exactly the kind of
  peeking that would be illegitimate with p-values.
\end{frame}

\begin{frame}{Optional continuation vs.\ optional stopping: who's continuing?}
  \begin{itemize}
    \item This setting is intimately related to sequential testing /
    optional stopping (Chapter 7), but is \emph{motivated differently}.
    \item Here, the scientist who continues the study may be a
    \emph{different} person from the one who collected the first batch.
    Data collection may be expensive -- the first scientist may simply have
    run out of money, or given up because the evidence looked unpromising.
    \item We call this setting \emph{optional continuation}, in contrast to
    \emph{optional stopping} (Chapter 7), where a \emph{single} experimenter
    controls if/when to stop a single ongoing experiment.
    \item Mathematically the multiplicative combination rule is identical;
    the distinction is about \emph{who} is making the post-hoc decision, and
    why.
  \end{itemize}
\end{frame}

\begin{frame}{The ``living meta-analysis'' story}
  \textbf{Scenario.} A drug (e.g.\ a pandemic vaccine) is expensive to test.
  \begin{itemize}
    \item Country A runs an expensive trial; the findings are impressive but
    Country B still deems it too risky to launch without its own evidence.
    \item Country B, piggy-backing on Country A's success, runs its own
    trial and \emph{appends} its e-value evidence to Country A's.
    \item Country C, who wasn't even considering the drug, may now be
    convinced to run its own trial too.
  \end{itemize}
  As long as every country's evidence is summarized as an e-value, the
  overall product $E^{(1)}E^{(2)}E^{(3)}\cdots$ is easy to update and
  monitor sequentially across this unplanned, ad hoc sequence of trials.
  \textbf{Toy numbers:} wealthier countries might launch a vaccine program
  once the combined e-value exceeds $20$ (i.e.\ post-hoc level $\le0.05$),
  while poorer, more risk-averse countries wait until it exceeds $100$
  (post-hoc level $\le0.01$) before acting \emph{without even running their
  own trial}.
\end{frame}

\begin{frame}[allowframebreaks]{Three important caveats}
  \begin{block}{1. No hiding unfavorable findings}
    No participant may discard or hide their result, even if their own batch
    gives a small e-value (less than $1$, i.e.\ evidence \emph{for} $H_0$).
    Once the decision to run the experiment is made, the resulting e-value
    \emph{must} be incorporated into the combined evidence.
  \end{block}

  \begin{block}{2. The ``poison pill'' problem, and two fixes}
    A single scientist reporting an e-value of exactly $0$ can permanently
    destroy \emph{all} previously accumulated evidence (since the running
    product becomes $0$ forever after). Two ways to guard against this:
    \begin{itemize}
      \item Agree in advance to use e-values that are never exactly $0$ by
      design (e.g.\ log-optimal e-values);
      \item Or agree to combine via
      $\prod_{i=1}^n\big(\tfrac{1+E_i}{2}\big)$, so no single batch can
      destroy more than half of the existing evidence -- generalizable to
      $\prod_{i=1}^n(\lambda_i+(1-\lambda_i)E_i)$ with $\lambda_i\in[0,1]$
      \emph{fixed before seeing $E_i$}.
    \end{itemize}
    Either way: the $i$-th e-value's construction, and how it is combined
    with the previous ones, must be specified \emph{before} seeing the
    $i$-th batch of data.
  \end{block}

  \framebreak
  \begin{block}{3. The decision to continue cannot be based on the p-value ``not being small enough''}
    The decision to continue collecting data at a second stage \emph{cannot}
    be triggered by ``the first p-value wasn't quite below $0.05$'' --
    testing at the first stage with that p-value already used up the entire
    type-I error budget. One must commit either to using a p-value with a
    single, fixed batch, \emph{or} to using an e-value with the possibility
    of extending the analysis -- not switch between the two after the fact.
  \end{block}
\end{frame}

\begin{frame}{Why can't we just combine p-values the same way?}
  It is tempting to think the second batch (say Country B's trial data) is
  ``independent'' of the first, and to merge the two p-values with a
  standard rule for independent p-values (e.g.\ Fisher's method).
  \textbf{This would be a mistake.}
  \begin{itemize}
    \item The second p-value is \emph{not} really independent of the first,
    despite being computed on ``separate'' data: it would not even exist had
    the first p-value not looked promising enough to justify continuing!
    \item So the \emph{number} of p-values collected is itself random, and
    the p-values are \emph{sequentially dependent} in a complicated way that
    standard combination rules do not account for.
    \item Such post-hoc decisions to run further trials, or to fold in
    accumulated data from separate trials, are a defining feature of
    e-value-based evidence collection -- and a fundamental obstacle for
    p-value-based analysis.
  \end{itemize}
\end{frame}

\begin{frame}[fragile]{Python: e-process vs.\ naive p-value peeking (\S4.3)}
\begin{lstlisting}
import numpy as np
from scipy.special import erfc
rng = np.random.default_rng(2)

theta0, theta1 = 0.5, 0.6
n_batches, batch_size, n_runs, alpha = 20, 10, 20_000, 0.05
threshold = 1/alpha                     # = 20

heads = rng.binomial(batch_size, theta0, size=(n_runs, n_batches))
tails = batch_size - heads

e_batch = (theta1/theta0)**heads * ((1-theta1)/(1-theta0))**tails
z = (heads - batch_size*theta0)/np.sqrt(batch_size*theta0*(1-theta0))
p_batch = np.clip(0.5*erfc(z/np.sqrt(2)), 1e-300, 1.0)

e_process = np.cumprod(e_batch, axis=1)          # the e-process
ever_reject_e = (np.cumsum(e_process >= threshold, axis=1) > 0)
ever_reject_p = (np.cumsum(p_batch < alpha, axis=1) > 0)   # naive peeking

print("false-rejection rate, e-process (Ville):", ever_reject_e[:, -1].mean())
print("false-rejection rate, naive p peeking  :", ever_reject_p[:, -1].mean())
\end{lstlisting}
  \textbf{Output (20 looks, $\alpha=0.05$):} e-process monitoring stays at
  $\approx 0.028\le 0.05$ (Ville's inequality holds), while naive repeated
  p-value peeking inflates to $\approx 0.675$ -- more than \emph{thirteen
  times} the nominal level. See \texttt{figures/fig\_optional\_continuation.pdf}.
\end{frame}

\begin{frame}{Section 4.3 take-aways}
  \begin{itemize}
    \item E-values multiply: fresh e-values on fresh (even post-hoc decided)
    batches combine into a valid e-value for the whole history, by the
    tower property $\mathbb E^{\mathbb P}[E_2\mid E_1]=1$.
    \item This licenses \emph{optional continuation}: deciding whether, and
    how, to gather more data \emph{after} seeing preliminary results ---
    something ordinary p-values cannot support once the first test's error
    budget is spent.
    \item Three caveats keep the guarantee honest: no hiding results, guard
    against a poison-pill e-value of $0$, and never let the p-value's own
    "closeness to significance" trigger continuation.
    \item Sequential monitoring of the e-process is controlled by Ville's
    inequality (previewed here, developed fully in Chapter 7).
  \end{itemize}
\end{frame}

% =======================================================================
\section{Incentive-Compatible Principal-Agent Hypothesis Testing}

\begin{frame}{Section 4.4}
  \tableofcontents[currentsection]
\end{frame}

\begin{frame}{Intuition: the regulator and the pharma company}
  \begin{itemize}
    \item Two parties: a \emph{principal} (e.g.\ a regulator granting drug
    licenses) and an \emph{agent} (e.g.\ a pharmaceutical company running a
    trial).
    \item Two features make this different from ordinary testing:
    \begin{enumerate}
      \item \textbf{Information asymmetry:} the agent may know far more
      about whether their drug actually works than the regulator does.
      \item \textbf{Conflicting incentives:} the regulator wants to keep
      ineffective drugs off the market; the agent wants a license to sell
      and profit, whether or not the drug works.
      \end{enumerate}
    \item Standard practice: if $P<0.05$, grant the license. We will see
    \emph{numerically} why this is not \emph{incentive compatible}, and how
    e-values fix it.
  \end{itemize}
\end{frame}

\begin{frame}{The incentive-compatibility problem: a numeric payoff table}
  A testing paradigm is \emph{incentive compatible} if it discourages
  \emph{malicious} agents (who know $H_0$ is true, e.g.\ they have a
  placebo) from even attempting a trial. Suppose a trial costs
  \$$10$ million, and the standard rule is: license granted iff $P<0.05$.

  \begin{block}{Toy payoff table for a malicious agent (drug is a placebo)}
    \begin{center}
    \begin{tabular}{lccc}
      \toprule
      Anticipated profit if approved & Trial cost & $\Pr(\text{approval}\mid H_0)$ & Expected net profit\\
      \midrule
      \$100 million & \$10 million & $0.05$ & $100(0.05)-10 = -\$5$M\\
      \$1{,}000 million (\$1B) & \$10 million & $0.05$ & $1000(0.05)-10=+\$40$M\\
      \bottomrule
    \end{tabular}
    \end{center}
  \end{block}
  \textbf{Reading the table:} if the anticipated profit is modest (\$100M),
  the malicious agent expects to \emph{lose} money and won't run the trial
  -- good. But if the anticipated profit is large enough (\$1B), the
  malicious agent's expected profit is strongly \emph{positive}: they will
  run the trial \emph{purely to gamble on the $5\%$ chance of approval}.
  The threshold rule is \emph{not} incentive compatible once profits are
  large enough.
\end{frame}

\begin{frame}{Why this matters at scale}
  \begin{block}{A quick back-of-envelope calculation}
    Suppose there are $20$ times as many ineffective drugs (placebos, in
    development) as truly effective ones, and $5\%$ of ineffective drugs get
    approved (the false-positive rate $\alpha=0.05$).
    \begin{center}
    \begin{tabular}{lccc}
      \toprule
      Drug type & Relative count & Approval rate & Expected \# approved\\
      \midrule
      Ineffective & $20$ & $5\%$ & $20\times0.05 = 1.0$\\
      Effective & $1$ & $\le 100\%$ (power) & $\le 1.0$\\
      \bottomrule
    \end{tabular}
    \end{center}
  \end{block}
  Since the expected number of \emph{ineffective} drugs approved ($1.0$) is
  already at least as large as the expected number of \emph{effective}
  drugs approved ($\le1.0$), \textbf{at least half of all approved drugs
  will be ineffective} -- a stark illustration of why the naive
  threshold-on-a-p-value protocol is dangerous at a population level, quite
  apart from the incentive problem above.
\end{frame}

\begin{frame}{Designing a better protocol}
  Another problem with the naive protocol: a malicious agent's potential
  profit is limited only by \emph{market size}, not by the strength of their
  statistical evidence. We would like a protocol with two properties:
  \begin{enumerate}
    \item \textbf{Incentive compatible:} malicious agents have no incentive
    to participate, regardless of market size.
    \item \textbf{Evidence-scaled reward:} if an agent participates, their
    permitted sales/profit should \emph{increase with the strength of the
    statistical evidence} for effectiveness.
  \end{enumerate}
  We will see that both are easily achieved with a \emph{statistical
  contract} built from e-values.
\end{frame}

\begin{frame}{The statistical contract}
  Let $\theta\in\Theta$ be the drug's true quality (known to the agent),
  with $\Theta_0\subsetneq\Theta$ the ``placebo'' parameter values the
  regulator wants to keep off the market.

  \begin{block}{Protocol: the agent follows these steps, in order}
    \begin{enumerate}
      \item Choose a nonnegative \emph{license function} $f$ from a menu of
      options $\mathcal F$ (offered by the regulator).
      \item Incur a fixed cost $C$ to run a trial, collecting data
      $Z\overset{\mathrm d}{\sim}\mathbb P_\theta$.
      \item Receive a license to profit at most $L=f(Z)$.
    \end{enumerate}
  \end{block}
  The agent's expected profit is $\mathbb E^{\mathbb P_\theta}[f(Z)]-C$; they
  participate iff this is positive. The regulator wants this to be
  \emph{negative} whenever $\theta\in\Theta_0$ (a placebo).
\end{frame}

\begin{frame}{Proposition 4.11: characterizing incentive compatibility}
  \begin{block}{Proposition 4.11}
    The statistical contract is incentive compatible if and only if
    $f(Z)/C$ is an e-variable for $\{\mathbb P_\theta:\theta\in\Theta_0\}$,
    for every $f\in\mathcal F$.
  \end{block}
  \textbf{Proof.} Incentive compatible means
  $\mathbb E^{\mathbb P_\theta}[f(Z)]\le C$ for all $\theta\in\Theta_0$,
  i.e.\ $\mathbb E^{\mathbb P_\theta}[f(Z)/C]\le 1$ for all
  $\theta\in\Theta_0$ -- precisely the definition of $f(Z)/C$ being an
  e-variable for the placebo class. $\blacksquare$

  \medskip
  \textbf{This is exactly Section 4.2's decision-maker framework}, with the
  regulator constructing a menu of \emph{scaled e-values} as the allowed
  license functions.
\end{frame}

\begin{frame}{A p-to-e calibrator fixes the naive protocol}
  Suppose evidence is summarized as a p-value $P$. It might be tempting to
  promise a license with profit margin \emph{inversely proportional} to
  $P$, i.e.\ $f(Z)=C/P$.
  \begin{itemize}
    \item This is \textbf{not} incentive compatible: since
    $\mathbb E^{\mathbb P_\theta}[1/P]=\infty$ (as we computed in Section 4.1!),
    the expected profit under the null is \emph{infinite}.
    \item Instead, awarding a license of
    \[
      f(Z) = \frac{C}{2\sqrt{P}}
    \]
    \emph{is} incentive compatible: $p\mapsto 1/(2\sqrt p)$ is a
    \emph{p-to-e calibrator}, i.e.\ $\mathbb E^{\mathbb P_\theta}[1/(2\sqrt P)]\le1$
    under $H_0$, so $\mathbb E^{\mathbb P_\theta}[f(Z)]\le C$ for all
    $\theta\in\Theta_0$ as required.
    \item Intuition: the calibrator ``tempers'' the explosive $1/P$ tail
    just enough to keep its expectation finite -- turning a p-value into a
    genuine e-value.
  \end{itemize}
\end{frame}

\begin{frame}[fragile]{Python: checking incentive compatibility of license functions (\S4.4)}
\begin{lstlisting}
import numpy as np
rng = np.random.default_rng(3)

M = 2_000_000
P = rng.uniform(0, 1, size=M)     # p-value under a placebo (H0 true)
C = 10.0                          # trial cost, in $ millions

f_naive = C / P                   # naive: profit inversely prop. to P
f_calibrated = C / (2*np.sqrt(P)) # p-to-e calibrator: 1/(2 sqrt(p))

print("E[f_naive(Z)]/C       :", np.mean(f_naive)/C, " (theory: infinite)")
print("E[f_calibrated(Z)]/C  :", np.mean(f_calibrated)/C, " (theory: <= 1)")
\end{lstlisting}
  Empirically, \texttt{f\_naive} keeps growing without bound as $M$
  increases (occasional tiny $P$ dominate the average), confirming it is
  \emph{not} incentive compatible; \texttt{f\_calibrated} stabilizes at a
  value $\le 1$, confirming $f(Z)/C$ is a genuine e-variable.
\end{frame}

\begin{frame}{Connecting back to Section 4.2: two decision makers}
  \begin{enumerate}
    \item \textbf{The regulator moves first:} they have no objective to
    maximize, but must meet the constraint of incentive compatibility. This
    is achieved by offering a menu $\mathcal F$ composed of (scaled)
    e-values; to encourage effective drugs to participate, $\mathcal F$
    should otherwise be \emph{unrestricted} -- the set of \emph{all} scaled
    e-values.
    \item \textbf{The agent moves second:} facing a data-dependent loss
    (their payout depends on the data they will collect), and with
    incentive compatibility already guaranteed by the regulator's choice of
    $\mathcal F$, their decision is simple: participate iff some $f\in\mathcal F$
    is profitable in expectation under their true (privately known)
    $\theta$. Malicious agents (who know $\theta\in\Theta_0$) never will;
    agents with genuinely effective drugs may.
  \end{enumerate}
  \textbf{Summary:} smart use of data-dependent loss functions, built from a
  menu of scaled e-values, considerably simplifies the decision problem for
  \emph{both} parties.
\end{frame}

\begin{frame}{Section 4.4 take-aways}
  \begin{itemize}
    \item Fixed p-value thresholds are not incentive compatible: a
    malicious agent with a large enough potential market will always find
    it profitable to gamble on the false-positive rate.
    \item Incentive compatibility is exactly characterized (Prop.\ 4.11) by
    requiring the license/profit function, scaled by the trial cost, to be
    an e-variable for the ``placebo'' null class.
    \item A p-to-e calibrator (e.g.\ $1/(2\sqrt P)$) converts a naive,
    incentive-\emph{in}compatible p-value-based reward into an incentive
    compatible one.
    \item This is the same decision-theoretic machinery as Section 4.2,
    applied to a two-player (principal-agent) game instead of a single
    decision maker.
  \end{itemize}
\end{frame}

% =======================================================================
\section*{Chapter Summary}

\begin{frame}{Chapter 4: the big picture}
  \begin{itemize}
    \item \textbf{4.1} Post-hoc validity: e-values (uniquely) support
    choosing your significance level \emph{after} seeing the data, with a
    guarantee that survives even adversarial choices of $\hat\alpha$.
    Ordinary p-values do not: $\mathbb E[1/P]=\infty$.
    \item \textbf{4.2} Post-hoc decisions: e-value-based decision rules
    satisfy risk constraints even for non-binary decisions and
    data-dependent decision problems; p-value-based rules can silently
    violate the risk budget (Example 4.10).
    \item \textbf{4.3} Optional continuation: e-values combine
    multiplicatively across post-hoc-decided follow-up experiments,
    licensing ``living'' analyses that grow evidence over time; p-values
    cannot be safely recombined this way.
    \item \textbf{4.4} Principal-agent testing: incentive compatibility of a
    statistical contract is exactly captured by requiring scaled license
    functions to be e-variables -- protecting a regulator from agents who
    already know the truth.
    \item \textbf{Common thread:} in all four settings, the object that
    remains valid under adversarial, data-dependent manipulation is always
    the e-value.
  \end{itemize}
\end{frame}

\begin{frame}{Bibliographical note}
  \begin{itemize}
    \item \S4.1 combines results explicit or implicit in Wang and Ramdas
    (2022), Xu et al.\ (2024), and Koning (2023a), in the context of
    post-hoc $\alpha$ for testing or FDR/FCR control. The sufficient
    condition in Remark 4.2 is due to Blanchard and Roquain (2008).
    \item \S4.2 is based on Grünwald (2024), with some presentation
    differences.
    \item Grünwald et al.\ (2024a) coined the term \emph{optional
    continuation} (\S4.3); further explored by Shafer (2023) and Ter Schure
    and Grünwald (2022) (application to meta-analysis).
    \item The incentive-compatible principal-agent problem of \S4.4 was
    formulated and studied by Bates et al.\ (2022).
  \end{itemize}
\end{frame}

\end{document}
